% Switch footnotes to symbols for title page only
\renewcommand{\thefootnote}{\fnsymbol{footnote}}

\thispagestyle{empty}

\begin{center}
\vspace*{3cm}

{\Large Internalizing Environmental Risk: Insurance Design and Firm Behavior in Hazardous Industries}

\vspace{1cm}

{\normalsize \href{https://www.kalebkjavier.com}{Kaleb K. Javier}\footnote{University of California, Berkeley, Department of Agricultural and Resource Economics. Email: \texttt{kalebkja@berkeley.edu}. Website: \texttt{kalebkjavier.com}. I thank Meredith Fowlie, Billy Pizer, James Sallee, Nano Barahona, Ben Handel, Matt Backus, and Cailin Slattery for their invaluable comments and feedback. I also thank the participants and discussants at Camp Resources XXXI and the AERE Summer Conference 2025 for helpful discussions. All remaining errors are my own.}}

\vspace{0.2cm}

{\normalsize March 16, 2026}

\end{center}

\vspace{.2cm}

\begin{center}
{\small \textbf{Abstract}}
\end{center}

\begin{quote}
{\small \noindent Public trust funds in thirty-eight U.S. states insure underground storage tank
(UST) leaks at a uniform flat fee, even though leak risk varies widely with observable
traits such as tank age and wall construction. This paper quantifies how replacing
flat-fee pooling with actuarially priced coverage changes firm behavior and environmental
outcomes. I exploit Texas's 1999 closure of its trust fund, which forced all owners into
a private market that sets premiums according to facility risk. A facility-year panel
covering 1990--2021 links EPA inventories with administrative microdata for Texas and
eighteen control states that retained public funds. I construct a make-model comparison
sample that matches Texas facilities to control-state counterparts on the observable risk
characteristics, tank age bin and wall construction type, that private insurers price,
ensuring the comparison group is balanced on the dimensions along which the treatment
effect is theoretically heterogeneous. Difference-in-differences and event-study estimates
show that risk-rated premiums increased annual closure rates by two to three percentage
points, shifted closures from market exit to on-site replacement with safer equipment,
and reduced the probability of a reported leak by about one percentage point, nearly the
entire pre-policy baseline for the high-risk single-wall cohort. The treatment effect is
concentrated among older and single-walled facilities, consistent with the
premium-gradient mechanism, and is near zero for young double-walled tanks. The findings
provide the first causal micro-evidence that cost-based environmental liability insurance
can internalize pollution risk more effectively than uniform public pooling, offering
clear guidance for states that are reconsidering the structure of UST insurance programs.}
\end{quote}

% Switch back to arabic numerals for the rest of the document
\renewcommand{\thefootnote}{\arabic{footnote}}
\setcounter{footnote}{0}

\newpage